\chapter{Trabalho de uma força}\label{ch:g11:work-of-force}

Empurre um guarda-roupa pela sala e você trabalhou; empurre uma parede a
tarde inteira e, digam o que disserem os seus músculos, a física conta
zero. O \omterm{def:g11:work-of-force:work}{trabalho} é o que uma \omterm{def:g10:inertia:force}{força} entrega \emph{ao longo de um
movimento} --- \omterm{def:g10:energy-conservation:energy}{energia} em trânsito
(\cref{ch:g10:energy-conservation}), paga em \omterm{def:g11:work-of-force:work}{joules}. Este capítulo dá
preço às \omterm{def:g10:inertia:force}{forças} do \cref{ch:g11:forces-and-motion} --- puxões em ângulo,
\omterm{def:g10:universal-gravitation:weight}{peso}, \omterm{def:g10:inertia:inventory}{atrito} --- e cronometra o ritmo do pagamento: a \omterm{def:g11:work-of-force:power}{potência}.

\section{Trabalho de uma força constante}

\begin{definition}[Trabalho de uma força constante]\label{def:g11:work-of-force:work}
Uma \omterm{def:g10:inertia:force}{força} constante $\vect F$ aplicada a um objeto que se move ao longo
de um segmento retilíneo de $A$ até $B$ realiza o
\emph{trabalho}\index{trabalho (de uma força)}
\[
W_{AB}(\vect F) = F \times AB \times \cos\theta ,
\]
em que $F$ é o módulo da \omterm{def:g10:inertia:force}{força} (em \unit{N}), $AB$ é o comprimento do
deslocamento (em \unit{m}) e $\theta$ é o ângulo entre $\vect F$ e
$\vect{AB}$. O trabalho se mede em \emph{joules}\index{joule}
(\unit{J}): um joule é o trabalho de uma \omterm{def:g10:inertia:force}{força} de um \omterm{def:g10:inertia:force}{newton} ao longo de
um \omterm{def:g10:orders-of-magnitude:unit}{metro} na sua própria direção.
\end{definition}

\begin{remark}\label{rem:g11:work-of-force:scal}
$F \times AB \times \cos\theta$ é exatamente o produto escalar
$\vect F \cdot \vect{AB}$ do volume de matemática deste ano. Só a
componente $F\cos\theta$ \emph{ao longo} do movimento rende alguma coisa.
\end{remark}

\begin{omfigure}
\begin{tikzpicture}[font=\small]
  \draw (-0.4,0) -- (6.4,0);
  \draw[fill=omExo!15] (0.4,0) rectangle (1.5,0.7);
  \coordinate (P) at (0.95,0.35);
  \draw[dashed, omExo] ($(P)+(2.29,1.61)$) -- ($(P)+(2.29,0)$);
  \draw[-Stealth, omDef, very thick] (P) -- ($(P)+(35:2.8)$)
    node[above] {$\vect F$};
  \draw[-Stealth, omMeth, thick] (P) -- ++(2.29,0)
    node[pos=0.75, below] {$F\cos\theta$};
  \draw (P) ++(1.1,0) arc[start angle=0, end angle=35, radius=1.1];
  \node at ($(P)+(17:1.38)$) {$\theta$};
  \draw[-Stealth, omThm, thick] (0.95,-0.7) -- (4.95,-0.7)
    node[midway, below] {$\vect{AB}$};
  \fill (0.95,-0.7) circle (1.2pt) node[below left] {$A$};
  \fill (4.95,-0.7) circle (1.2pt) node[below right] {$B$};
\end{tikzpicture}

{\small Só a projeção $F\cos\theta$ da \omterm{def:g10:inertia:force}{força} sobre o deslocamento realiza
\omterm{def:g11:work-of-force:work}{trabalho}: $W_{AB}(\vect F) = F \times AB \times \cos\theta$.}
\end{omfigure}

\begin{example}[Mala na correia]\label{ex:g11:work-of-force:suitcase}
Um viajante arrasta uma mala por \qty{300}{m} num terminal, puxando a
alça com $F = \qty{40}{N}$ a $\theta = \qty{50}{\degree}$ do piso:
$W = 40 \times 300 \times \cos\qty{50}{\degree} \approx \qty{7.7}{kJ}$.
Puxada na horizontal ($\theta = 0$), a mesma \omterm{def:g10:inertia:force}{força} entregaria
\qty{12}{kJ}: um terço do puxão não levanta nada.
\end{example}

\begin{definition}[Trabalho motor, resistente e nulo]\label{def:g11:work-of-force:sign}
O sinal de $\cos\theta$ classifica toda \omterm{def:g10:inertia:force}{força} em três regimes:
\begin{itemize}
  \item $\theta < \qty{90}{\degree}$: $W > 0$, o \omterm{def:g11:work-of-force:work}{trabalho} é
        \emph{motor}\index{trabalho motor} --- a \omterm{def:g10:inertia:force}{força} injeta \omterm{def:g10:energy-conservation:energy}{energia} no
        movimento;
  \item $\theta = \qty{90}{\degree}$: $W = 0$, a \omterm{def:g10:inertia:force}{força} não realiza
        \omterm{def:g11:work-of-force:work}{trabalho} algum, por maior que seja;
  \item $\theta > \qty{90}{\degree}$: $W < 0$, o \omterm{def:g11:work-of-force:work}{trabalho} é
        \emph{resistente}\index{trabalho resistente} --- a \omterm{def:g10:inertia:force}{força} drena
        \omterm{def:g10:energy-conservation:energy}{energia}. O caso puro é um freio, a $\qty{180}{\degree}$:
        \qty{800}{N} ao longo de \qty{45}{m} retiram
        $800 \times 45 \times (-1) = \qty{-36}{kJ}$ de um carro, e os
        despejam em discos quentes.
\end{itemize}
\end{definition}

\begin{omfigure}
\begin{tikzpicture}[font=\small]
  \draw (0,0) -- (2.6,0); \draw[fill=omExo!15] (0.9,0) rectangle (1.7,0.55);
  \draw[-Stealth, omDef, very thick] (1.3,0.28) -- ++(40:1.1)
    node[above] {$\vect F$};
  \draw[-Stealth, omThm, thick] (0.5,-0.5) -- (2.1,-0.5);
  \node at (1.3,2.0) {$W > 0$};
\end{tikzpicture}\qquad
\begin{tikzpicture}[font=\small]
  \draw (0,0) -- (2.6,0); \draw[fill=omExo!15] (0.9,0) rectangle (1.7,0.55);
  \draw[-Stealth, omDef, very thick] (1.3,0.55) -- ++(90:1.1)
    node[above] {$\vect F$};
  \draw[-Stealth, omThm, thick] (0.5,-0.5) -- (2.1,-0.5);
  \node at (1.3,2.0) {$W = 0$};
\end{tikzpicture}\qquad
\begin{tikzpicture}[font=\small]
  \draw (0,0) -- (2.6,0); \draw[fill=omExo!15] (0.9,0) rectangle (1.7,0.55);
  \draw[-Stealth, omDef, very thick] (1.3,0.28) -- ++(140:1.1)
    node[above] {$\vect F$};
  \draw[-Stealth, omThm, thick] (0.5,-0.5) -- (2.1,-0.5);
  \node at (1.3,2.0) {$W < 0$};
\end{tikzpicture}

{\small Os três regimes --- \omterm{def:g11:work-of-force:sign}{motor}, nulo e \omterm{def:g11:work-of-force:sign}{resistente} --- pelo ângulo
entre a \omterm{def:g10:inertia:force}{força} e o deslocamento (flecha vermelha).}
\end{omfigure}

\section{O trabalho do peso}

\begin{proposition}[Trabalho do peso]\label{prop:g11:work-of-force:weight}
Quando um objeto de massa $m$ vai de $A$ (altitude $z_A$) até $B$
(altitude $z_B$) por \emph{qualquer} caminho feito de segmentos
retilíneos, o seu \omterm{def:g10:universal-gravitation:weight}{peso} $\vect P$ realiza o \omterm{def:g11:work-of-force:work}{trabalho}
\[
W_{AB}(\vect P) = mg\,(z_A - z_B) = \pm\, mgh ,
\]
em que $h = \abs{z_A - z_B}$ é o desnível: $+mgh$ na descida e $-mgh$ na
subida --- só o desnível importa, e não o percurso.
\end{proposition}

\begin{proof}
Num único segmento retilíneo $AB$, o \omterm{def:g10:universal-gravitation:weight}{peso} (módulo $mg$, dirigido para
baixo) faz um ângulo $\theta$ com $\vect{AB}$, e $AB\cos\theta$ é
justamente a projeção do deslocamento sobre a vertical descendente:
$z_A - z_B$. Logo, $W = mg \times AB\cos\theta = mg\,(z_A - z_B)$. Numa
sequência de segmentos os \omterm{def:g11:work-of-force:work}{trabalhos} se somam e os desníveis se
telescopam: $mg(z_A - z_C) + mg(z_C - z_B) = mg(z_A - z_B)$.
\end{proof}

\begin{proposition}[Qualquer caminho]\label{prop:g11:work-of-force:anypath}
A mesma fórmula vale ao longo de qualquer caminho curvo de $A$ até $B$.
\end{proposition}
\admitted

\begin{remark}\label{rem:g11:work-of-force:anypath}
É plausível --- uma curva está tão perto quanto queiramos de uma linha
quebrada de muitos segmentos curtos ---, mas tornar honesto esse ``tão
perto quanto queiramos'' é cálculo, feito no
\cref{ch:g12:work-and-energy}.
\end{remark}

\begin{omfigure}
\begin{tikzpicture}[font=\small]
  \coordinate (A) at (0.7,2.6); \coordinate (B) at (4.4,0);
  \draw[dashed, omExo] (0,2.6) -- (5.2,2.6);
  \draw[dashed, omExo] (0,0) -- (5.2,0);
  \draw[Stealth-Stealth, omExo] (5.0,0.05) -- (5.0,2.55)
    node[midway, right] {$h$};
  \draw[omDef, very thick] (A) -- (B);
  \draw[omMeth, very thick] (A) -- (1.5,2.0) -- (3.1,2.0) --
    (2.3,0.9) -- (3.6,0.55) -- (B);
  \draw[-Stealth, omThm, thick] (2.55,1.85) -- ++(0,-0.8)
    node[below] {$\vect P$};
  \fill (A) circle (1.5pt) node[above] {$A$};
  \fill (B) circle (1.5pt) node[below] {$B$};
\end{tikzpicture}

{\small Reto (azul) ou sinuoso (verde), todo caminho de $A$ até $B$ desce
o mesmo $h$: o \omterm{def:g10:universal-gravitation:weight}{peso} realiza $+mgh$ nos dois.}
\end{omfigure}

\begin{example}[Caminhante]\label{ex:g11:work-of-force:hiker}
Um caminhante de \qty{78}{kg} (mochila incluída) sobe \qty{450}{m}, por
qualquer trilha: $W(\vect P) = -78 \times 9.81 \times 450 \approx
\qty{-344}{kJ}$, e exatamente $+\qty{344}{kJ}$ na descida, sejam quais
forem os desvios.
\end{example}

\section{O trabalho do atrito}

\begin{proposition}[Trabalho do atrito de deslizamento]\label{prop:g11:work-of-force:friction}
Uma \omterm{def:g10:inertia:force}{força} de \omterm{def:g10:inertia:inventory}{atrito} de módulo constante $f$, a cada instante oposta ao
movimento, realiza o \omterm{def:g11:work-of-force:work}{trabalho} $W = -fL$ ao longo de um caminho de
comprimento total $L$: sempre \omterm{def:g11:work-of-force:sign}{resistente} e proporcional ao
\emph{comprimento do caminho percorrido} --- e não ao deslocamento.
\end{proposition}

\begin{proof}
Em cada pedacinho retilíneo a \omterm{def:g10:inertia:force}{força} está a $\qty{180}{\degree}$ do
movimento, contribuindo com $-f \times (\text{o comprimento dele})$; os
pedaços somam $-fL$. Um caminho curvo é picado em pedacinhos do mesmo
jeito, com o limite admitido como no caso do \omterm{def:g10:universal-gravitation:weight}{peso}
(\cref{rem:g11:work-of-force:anypath}).
\end{proof}

\begin{remark}[Atrito contra peso: a grande divisão]\label{rem:g11:work-of-force:divide}
O \omterm{def:g11:work-of-force:work}{trabalho} do \omterm{def:g10:universal-gravitation:weight}{peso} esquece o caminho; o do \omterm{def:g10:inertia:inventory}{atrito} é um pedágio por \omterm{def:g10:orders-of-magnitude:unit}{metro}.
Numa ida e volta o \omterm{def:g10:universal-gravitation:weight}{peso} fecha em zero --- o que ele toma na subida
devolve na descida ---, enquanto o \omterm{def:g10:inertia:inventory}{atrito} cobra $-fL$ em cada trecho e
não devolve nada: o \omterm{def:g11:work-of-force:work}{trabalho} do \omterm{def:g10:universal-gravitation:weight}{peso} pode ser armazenado e recuperado
(\omterm{def:g10:energy-conservation:energy}{energia} potencial, \cref{ch:g11:mechanical-energy}), e o do \omterm{def:g10:inertia:inventory}{atrito} se
perde como calor.
\end{remark}

\begin{example}[Ida e volta]\label{ex:g11:work-of-force:roundtrip}
Um engradado desliza \qty{5.0}{m} pelo piso e \qty{5.0}{m} de volta,
contra $f = \qty{60}{N}$: o \omterm{def:g10:inertia:inventory}{atrito} realiza
$-60 \times 10 = \qty{-600}{J}$, embora o engradado termine onde
começou. \omterm{def:g10:universal-gravitation:weight}{Peso} e \omterm{def:g10:inertia:force}{força} normal, perpendiculares ao movimento: zero o tempo
todo.
\end{example}

\begin{remark}[E uma mola?]\label{rem:g11:work-of-force:spring}
Uma mola puxa com mais \omterm{def:g10:inertia:force}{força} quanto mais esticada estiver: a \omterm{def:g10:inertia:force}{força} dela
não é constante e a \cref{def:g11:work-of-force:work} não se aplica.
Calcular o \omterm{def:g11:work-of-force:work}{trabalho} dela significa cortar o alongamento em passos
infinitesimais e somar --- cálculo, entregue no
\cref{ch:g12:work-and-energy}.
\end{remark}

\section{Potência}

\begin{definition}[Potência]\label{def:g11:work-of-force:power}
A \emph{potência}\index{potência!mecânica} de uma \omterm{def:g10:inertia:force}{força} é o ritmo com que
ela realiza \omterm{def:g11:work-of-force:work}{trabalho}, $P = W/\Delta t$, em \emph{watts}\index{watt}
($\unit{W} = \unit{J/s}$) --- o mesmo watt do
\cref{ch:g11:circuits-and-power}. Uma \omterm{def:g10:orders-of-magnitude:unit}{unidade} anterior ao SI sobrevive
nas placas dos \omterm{def:g11:work-of-force:sign}{motores}: o \emph{cavalo-vapor}\index{cavalo-vapor},
$\qty{1}{hp} = \qty{736}{W}$.
\end{definition}

\begin{proposition}[Potência com velocidade constante]\label{prop:g11:work-of-force:pfv}
Se o objeto se move com velocidade escalar constante $v$, uma \omterm{def:g10:inertia:force}{força}
constante $\vect F$ a um ângulo $\theta$ do movimento entrega a \omterm{def:g11:work-of-force:power}{potência}
$P = F v \cos\theta$ --- em particular, $P = Fv$ para uma \omterm{def:g10:inertia:force}{força} ao longo
do movimento.
\end{proposition}

\begin{proof}
Num tempo $\Delta t$ o objeto percorre $d = v\,\Delta t$, de modo que
$W = F \times v\,\Delta t \times \cos\theta$; divida por $\Delta t$.
\end{proof}

\begin{example}[Guindaste]\label{ex:g11:work-of-force:crane}
Um guindaste iça um palete de \qty{600}{kg} a constantes
\qty{0.90}{m/s}: a \omterm{def:g10:inertia:force}{força} do cabo é igual ao \omterm{def:g10:universal-gravitation:weight}{peso},
$mg = \qty{5.89e3}{N}$, logo
$P = Fv = 5890 \times 0.90 \approx \qty{5.3}{kW} \approx \qty{7.2}{hp}$
--- uma parelha de sete cavalos dentro de uma caixa de aço.
\end{example}

\begin{method}[Subir com menos força: o ziguezague]\label{met:g11:work-of-force:zigzag}
Para elevar uma carga a uma altura $h$ com uma \omterm{def:g10:inertia:force}{força} limitada:
\begin{enumerate}
  \item A conta está fixada: pela
        \cref{prop:g11:work-of-force:weight}, qualquer rota até o topo
        custa o \omterm{def:g11:work-of-force:work}{trabalho} $mgh$.
  \item Estique o caminho: numa rampa de comprimento $L$ a \omterm{def:g10:inertia:force}{força}
        necessária (com velocidade constante, deixando o \omterm{def:g10:inertia:inventory}{atrito} de lado)
        vale $F = mgh/L$ --- dobre o comprimento e reduza a \omterm{def:g10:inertia:force}{força} à
        metade.
  \item Pague a sobretaxa: o \omterm{def:g10:inertia:inventory}{atrito} cobra por \omterm{def:g10:orders-of-magnitude:unit}{metro}
        (\cref{prop:g11:work-of-force:friction}), acrescentando $fL$.
\end{enumerate}
As estradas de montanha ziguezagueiam pelo passo 2: mesma altura, menos
\omterm{def:g10:inertia:force}{força}, mesmo \omterm{def:g11:work-of-force:work}{trabalho}.
\end{method}

\begin{omfigure}
\begin{tikzpicture}
\begin{axis}[omaxis, width=8.8cm, height=5.0cm,
  xlabel={comprimento do caminho $L$ (\unit{km})},
  ylabel={força necessária (\unit{kN})},
  xmin=0.8, xmax=6.2, ymin=0, ymax=3.2,
  domain=1.2:6, samples=100]
  \addplot[omDef, thick] {3.53/x};
  \addplot[only marks, mark=*, omThm] coordinates {(1.5,2.35) (5,0.71)};
  \node[font=\small, anchor=west] at (axis cs:1.65,2.35) {trilha reta};
  \node[font=\small, anchor=south west] at (axis cs:4.0,0.8) {estrada em ziguezague};
\end{axis}
\end{tikzpicture}

{\small Puxar \qty{1200}{kg} até $h = \qty{300}{m}$: o produto
$F \times L = mgh = \qty{3.53}{MJ}$ está fixado, de modo que a estrada
troca comprimento por \omterm{def:g10:inertia:force}{força} (veja o \cref{exo:g11:work-of-force:12}).}
\end{omfigure}

\section{Exercícios}

\begin{exercise}[$\star$]\label{exo:g11:work-of-force:1}
Uma \omterm{def:g10:inertia:force}{força} constante $F = \qty{40}{N}$ age ao longo de um deslocamento
retilíneo de \qty{50}{m}: calcule o seu \omterm{def:g11:work-of-force:work}{trabalho} para
$\theta = \qty{0}{\degree}$, $\qty{60}{\degree}$, $\qty{90}{\degree}$ e
$\qty{120}{\degree}$.
\end{exercise}

\begin{exercise}[$\star$]\label{exo:g11:work-of-force:2}
Uma caixa de ferramentas de \qty{12}{kg} é baixada \qty{2.5}{m} de uma
van e, mais tarde, içada de volta. Qual é o \omterm{def:g11:work-of-force:work}{trabalho} do \omterm{def:g10:universal-gravitation:weight}{peso} dela em cada
caso?
\end{exercise}

\begin{exercise}[$\star$]\label{exo:g11:work-of-force:3}
\omterm{def:g11:work-of-force:sign}{Motor}, \omterm{def:g11:work-of-force:sign}{resistente} ou nulo? (a) o \omterm{def:g10:universal-gravitation:weight}{peso} de uma maçã em queda; (b) o \omterm{def:g10:universal-gravitation:weight}{peso} de
uma bola em ascensão; (c) a \omterm{def:g10:inertia:force}{força} normal sobre uma caixa que desliza;
(d) o \omterm{def:g10:inertia:inventory}{atrito} sobre uma bicicleta que freia; (e) a tensão do fio sobre a
massa de um pêndulo que oscila.
\end{exercise}

\begin{exercise}[$\star$]\label{exo:g11:work-of-force:4}
Um guindaste levanta \qty{600}{kg} em \qty{25}{m} com velocidade
constante. Qual é o \omterm{def:g11:work-of-force:work}{trabalho} da \omterm{def:g10:inertia:force}{força} de içamento e o do \omterm{def:g10:universal-gravitation:weight}{peso}? Qual é a
soma deles, e por quê?
\end{exercise}

\begin{exercise}[$\star$]\label{exo:g11:work-of-force:5}
O \omterm{def:g11:work-of-force:sign}{motor} de uma escada rolante realiza \qty{90}{kJ} de \omterm{def:g11:work-of-force:work}{trabalho} a cada
minuto. Qual é a \omterm{def:g11:work-of-force:power}{potência} dele em \omterm{def:g11:work-of-force:power}{watts}? E em cavalos-vapor?
\end{exercise}

\begin{exercise}[$\star\star$]\label{exo:g11:work-of-force:6}
Um trenó é rebocado por \qty{200}{m} por uma corda de tensão \qty{120}{N}
a $\qty{30}{\degree}$ acima da neve; o \omterm{def:g10:inertia:inventory}{atrito} vale \qty{45}{N}. Qual é o
\omterm{def:g11:work-of-force:work}{trabalho} de cada uma das quatro \omterm{def:g10:inertia:force}{forças} (corda, \omterm{def:g10:inertia:inventory}{atrito}, \omterm{def:g10:universal-gravitation:weight}{peso}, normal) e o
total?
\end{exercise}

\begin{exercise}[$\star\star$]\label{exo:g11:work-of-force:7}
Um ciclista pedala a constantes \qty{25}{km/h} contra um arrasto total de
\qty{18}{N}. Qual é a \omterm{def:g11:work-of-force:power}{potência} entregue? E o \omterm{def:g11:work-of-force:work}{trabalho} realizado em uma
hora?
\end{exercise}

\begin{exercise}[$\star\star$]\label{exo:g11:work-of-force:8}
Um esquiador de \qty{65}{kg} desce \qty{120}{m} de altura, por uma pista
íngreme de \qty{800}{m} ou por uma trilha suave de \qty{1500}{m}; o
\omterm{def:g10:inertia:inventory}{atrito} nos esquis é de \qty{30}{N} nas duas. Qual é o \omterm{def:g11:work-of-force:work}{trabalho} do \omterm{def:g10:universal-gravitation:weight}{peso} e
o do \omterm{def:g10:inertia:inventory}{atrito} em cada rota? Que \omterm{def:g11:work-of-force:work}{trabalho} se importa com a rota?
\end{exercise}

\begin{exercise}[$\star\star$]\label{exo:g11:work-of-force:9}
Uma caixa de \qty{20}{kg} é empurrada \qty{4.0}{m} rampa acima, com
inclinação de $\qty{25}{\degree}$, com velocidade constante e contra
\qty{30}{N} de \omterm{def:g10:inertia:inventory}{atrito}. Calcule os \omterm{def:g11:work-of-force:work}{trabalhos} do \omterm{def:g10:universal-gravitation:weight}{peso} e do \omterm{def:g10:inertia:inventory}{atrito} e, em
seguida (a partir da velocidade constante), o \omterm{def:g11:work-of-force:work}{trabalho} e o módulo do
empurrão, paralelo à rampa.
\end{exercise}

\begin{exercise}[$\star\star$]\label{exo:g11:work-of-force:10}
Um carro segue a \qty{130}{km/h} com o \omterm{def:g11:work-of-force:sign}{motor} entregando \qty{50}{kW} às
rodas. Que \omterm{def:g10:inertia:force}{força} \omterm{def:g11:work-of-force:sign}{resistente} total ele enfrenta?
\end{exercise}

\begin{exercise}[$\star\star$]\label{exo:g11:work-of-force:11}
Um cavalo de canal reboca uma barcaça a \qty{1.0}{m/s}, com a corda de
tensão \qty{800}{N} fazendo $\qty{20}{\degree}$ com o caminho de sirga.
Calcule a \omterm{def:g11:work-of-force:power}{potência} do cavalo, em \omterm{def:g11:work-of-force:power}{watts} e em cavalos-vapor. Comente.
\end{exercise}

\begin{exercise}[$\star\star\star$]\label{exo:g11:work-of-force:12}
Um carro de \qty{1200}{kg} sobe $h = \qty{300}{m}$ até um passo, por uma
trilha reta de $20\%$ de inclinação ou por uma estrada em ziguezague de
$6\%$ (inclinação $= \sin\alpha$). Desprezando o \omterm{def:g10:inertia:inventory}{atrito}, com velocidade
constante, calcule para cada rota o comprimento, a \omterm{def:g10:inertia:force}{força} necessária ao
longo da rampa e o \omterm{def:g11:work-of-force:work}{trabalho} dela. Conclua em uma frase.
\end{exercise}

\begin{exercise}[$\star\star\star$]\label{exo:g11:work-of-force:13}
Um guincho ergue \qty{250}{kg} de telhas \qty{15}{m} por um andaime a
constantes \qty{0.40}{m/s}. Calcule o \omterm{def:g11:work-of-force:work}{trabalho} da \omterm{def:g10:inertia:force}{força} do cabo, a
duração, a \omterm{def:g11:work-of-force:power}{potência} mecânica (duas vezes: $W/\Delta t$ e $Fv$) e a
\omterm{prop:g11:circuits-and-power:power}{potência elétrica} consumida com $70\%$ de \omterm{def:g10:energy-conservation:efficiency}{rendimento} do \omterm{def:g11:work-of-force:sign}{motor}.
\end{exercise}

\begin{exercise}[$\star\star\star$]\label{exo:g11:work-of-force:14}
Um guarda-roupa é arrastado contra $f = \qty{140}{N}$ de \omterm{def:g10:inertia:inventory}{atrito} de um
canto da sala até o canto oposto: em linha reta pela diagonal de
\qty{6.0}{m}, ou \qty{3.6}{m} e depois \qty{4.8}{m} ao longo das
paredes. Qual é o \omterm{def:g11:work-of-force:work}{trabalho} do \omterm{def:g10:inertia:inventory}{atrito} em cada rota, e depois numa ida e
volta pela diagonal? O que o \omterm{def:g10:universal-gravitation:weight}{peso} faz nessa ida e volta, e que diferença
profunda entre as duas \omterm{def:g10:inertia:force}{forças} isso expõe?
\end{exercise}

\begin{exercise}[$\star\star\star$]\label{exo:g11:work-of-force:15}
Um ciclista (\qty{80}{kg} com a bicicleta) sobe um passo: \qty{600}{m} de
altura ao longo de \qty{12}{km} de estrada, a constantes \qty{15}{km/h},
contra \qty{14}{N} de arrasto e \omterm{def:g11:circuits-and-power:resistance}{resistência} de rolamento. Calcule o
\omterm{def:g11:work-of-force:work}{trabalho} fornecido contra a gravidade, contra o \omterm{def:g10:inertia:inventory}{atrito} e ao todo; a
duração da subida; e a \omterm{def:g11:work-of-force:power}{potência} média, em \omterm{def:g11:work-of-force:power}{watts} e em cavalos-vapor.
\end{exercise}

\section{Problema: o dia da mudança}

\begin{problem}[{Problema de fim de semana --- a física paga por joule:
um dia de arrastar, içar e subir curvas fechadas, terminando com a
humilhante descoberta do que uma chaleira acha do trabalho honesto}]
\label{pb:g11:work-of-force:1}
Uma equipe de mudanças esvazia um depósito, sobe a carga até um
apartamento de terceiro andar e depois leva a van até um povoado no alto
de um morro. Cada tarefa tem preço em \omterm{def:g11:work-of-force:work}{joules}; a chaleira da cozinha vai
auditar o dia. Caixas: \qty{30}{kg}.

\medskip
\noindent\textbf{Parte I --- Pelo piso do depósito.}
Uma caixa é empurrada com uma \omterm{def:g10:inertia:force}{força} horizontal $F = \qty{110}{N}$ por
$d = \qty{8.0}{m}$; o \omterm{def:g10:inertia:inventory}{atrito} de deslizamento vale $f = \qty{95}{N}$.
\begin{enumerate}
  \item Calcule o \omterm{def:g11:work-of-force:work}{trabalho} do empurrão.
  \item Calcule o \omterm{def:g11:work-of-force:work}{trabalho} do \omterm{def:g10:inertia:inventory}{atrito}, o do \omterm{def:g10:universal-gravitation:weight}{peso} e o da \omterm{def:g10:inertia:force}{força} normal
        (justifique os zeros).
  \item Qual é o \omterm{def:g11:work-of-force:work}{trabalho} total sobre a caixa e o significado do sinal
        dele (a caixa partiu do repouso)? Que empurrão mantém a
        velocidade constante (\cref{ch:g11:forces-and-motion})?
  \item O canto oposto pode ser alcançado pela diagonal de \qty{10.0}{m}
        ou por duas paredes (\qty{6.0}{m} e depois \qty{8.0}{m}). Qual é
        o \omterm{def:g11:work-of-force:work}{trabalho} do \omterm{def:g10:inertia:inventory}{atrito} em cada rota?
  \item Numa ida e volta pela diagonal, calcule o \omterm{def:g11:work-of-force:work}{trabalho} do \omterm{def:g10:inertia:inventory}{atrito} e o
        \omterm{def:g11:work-of-force:work}{trabalho} do \omterm{def:g10:universal-gravitation:weight}{peso}. Que \omterm{def:g10:inertia:force}{força} devolve?
\end{enumerate}

\medskip
\noindent\textbf{Parte II --- Até o terceiro andar ($h = \qty{9.0}{m}$).}
\begin{enumerate}[resume]
  \item Calcule o \omterm{def:g11:work-of-force:work}{trabalho} do \omterm{def:g10:universal-gravitation:weight}{peso} numa caixa erguida até o apartamento.
        Que \omterm{def:g11:work-of-force:work}{trabalho} a equipe precisa fornecer, no mínimo?
  \item Carregada escada acima por um carregador de \qty{70}{kg}, a
        gravidade cobra pelo carregador \emph{e} pela caixa: qual é o
        \omterm{def:g11:work-of-force:work}{trabalho} total contra a gravidade e a fração efetivamente gasta
        com a caixa?
  \item Uma rampa para carrinho tem \qty{30}{m} de comprimento para a
        mesma subida de \qty{9.0}{m}, com \qty{40}{N} de \omterm{def:g10:inertia:inventory}{atrito} de
        rolamento. Com velocidade constante, calcule o empurrão
        necessário ao longo da rampa e o \omterm{def:g11:work-of-force:work}{trabalho} dele.
  \item Compare a rampa com um içamento vertical direto: \omterm{def:g10:inertia:force}{força}
        necessária e \omterm{def:g11:work-of-force:work}{trabalho} realizado. O que a rampa compra, e a que
        preço?
  \item Um sistema de corda e roldana ergue a caixa a \qty{0.50}{m/s}.
        Qual é a duração e a \omterm{def:g11:work-of-force:power}{potência} mecânica?
  \item Recupere essa \omterm{def:g11:work-of-force:power}{potência} como $P = Fv$. O \omterm{def:g11:work-of-force:sign}{motor} do guincho tem
        $60\%$ de \omterm{def:g10:energy-conservation:efficiency}{rendimento}: qual é a \omterm{prop:g11:circuits-and-power:power}{potência elétrica} consumida?
\end{enumerate}

\medskip
\noindent\textbf{Parte III --- A estrada em ziguezague.}
A van carregada (\qty{3500}{kg}) precisa chegar a um povoado
\qty{240}{m} acima do vale: pela ladeira reta, \qty{1.0}{km} a $24\%$ de
inclinação, ou pela estrada asfaltada, \qty{8.0}{km} a $3.0\%$
(inclinação $= \sin\alpha$).
\begin{enumerate}[resume]
  \item Confira que as duas rotas sobem os mesmos \qty{240}{m}.
  \item Com velocidade constante e desprezando o \omterm{def:g10:inertia:inventory}{atrito}, calcule a \omterm{def:g10:inertia:force}{força}
        motriz necessária em cada rota.
  \item Qual é o \omterm{def:g11:work-of-force:work}{trabalho} da \omterm{def:g10:inertia:force}{força} motriz em cada rota? Compare com
        $mgh$.
  \item A \omterm{def:g11:circuits-and-power:resistance}{resistência} de rolamento vale \qty{400}{N} nas duas: calcule o
        \omterm{def:g11:work-of-force:work}{trabalho} extra em cada rota. Quanto custa o conforto do
        ziguezague?
  \item A \qty{36}{km/h}, calcule a \omterm{def:g11:work-of-force:power}{potência} de \omterm{def:g11:work-of-force:sign}{motor} necessária em cada
        rota (em \unit{kW} e em \unit{hp}). Por que as estradas
        ziguezagueiam?
\end{enumerate}

\medskip
\noindent\textbf{Parte IV --- O balanço do dia.}
\begin{enumerate}[resume]
  \item A equipe arrasta $60$ caixas pelo piso (com o empurrão da questão
        1) e iça as $60$ pelos \qty{9.0}{m} (questão 6). Qual é o
        \omterm{def:g11:work-of-force:work}{trabalho} mecânico total entregue às caixas?
  \item Acrescente o próprio corpo de quem sobe a escada: $20$ viagens de
        subida a \qty{70}{kg} (as devoluções da descida vão para joelhos
        doloridos, e não para o balanço). Qual é o total do dia?
  \item A equipe trabalhou $8.0$ horas: qual é a \omterm{def:g11:work-of-force:power}{potência} mecânica média?
        Quanto tempo a chaleira de \qty{2.0}{kW} levaria para gastar a
        mesma \omterm{def:g10:energy-conservation:energy}{energia}?
  \item Final: em duas frases, diga o que o \omterm{def:g11:work-of-force:work}{joule} paga --- e o que não
        paga (a parede empurrada a tarde inteira, a caixa segurada imóvel)
        --- e por que as equipes de mudança têm guinchos, rampas e
        estradas em ziguezague em vez de músculos maiores.
\end{enumerate}
\end{problem}
